\documentclass[11pt]{article}

% ----- Packages -----
\usepackage[margin=1in]{geometry}
\usepackage{amsmath, amssymb, amsthm}
\usepackage{braket}          % \ket, \bra, \braket, \ketbra
\usepackage{mathtools}
\usepackage{bm}
\usepackage{enumitem}
\usepackage{xcolor}
\usepackage{hyperref}
\hypersetup{colorlinks=true, linkcolor=blue!55!black, urlcolor=blue!55!black}

% ----- Convenience macros -----
\newcommand{\Tr}{\operatorname{Tr}}
\newcommand{\supp}{\operatorname{supp}}
\newcommand{\argmax}{\operatorname*{arg\,max}}
\DeclareMathOperator{\indic}{\mathbb{1}}
\newcommand{\R}{\mathbb{R}}
\newcommand{\onenorm}[1]{\left\lVert #1 \right\rVert_{1}}

\title{\textbf{Quantum Information Theory}\\[4pt]
Operational Distance for Quantum States:\\
Trace Distance \& Helstrom Discrimination}
\author{}
\date{Talk 2 --- June 25, 2026}

\begin{document}
\maketitle

\noindent\emph{Main idea:} introduce a distance measure between two quantum
states and relate it to something physically measurable.

%===============================================================
\section{Classical Case}
%===============================================================

We are given a sample realization $x$ (with $x\in\R$, say) that could be drawn
from one of two probability distributions, $p(x)$ or $q(x)$, with probability
$\tfrac12$ each.

\medskip
Given $\{x,\,p(x),\,q(x)\}$ we want to make the best choice / guess of which
pdf the sample $x$ came from. We want a \emph{decision rule}, optimal in some
sense, that chooses $X\sim p(x)$ or $X\sim q(x)$ with the highest probability
of success, $p^{\mathrm{succ}}$.

\[
\begin{array}{ccc}
X\sim p(x) \\[2pt]
\quad ? \searrow \\
& \boxed{\;\text{Decision Rule}\;}
& \longrightarrow\ \{P,\,Q\} \ \rightsquigarrow\ p^{\mathrm{succ}} \\
\quad ? \nearrow \\[2pt]
X\sim q(x)
\end{array}
\]

%---------------------------------------------------------------
\subsection{Systematically designing an optimal decision rule}
%---------------------------------------------------------------

Let $d$ be a random variable taking values in $\{P,Q\}$, representing the two
distributions $p(x),\,q(x)$ (a label). We observe a random variable $X$ drawn
from one of the two distributions:
\[
p(X\mid d=P) = p(X), \qquad
p(X\mid d=Q) = q(X).
\]
We have the (uniform) prior
\[
P(d=P) = P(d=Q) = \tfrac12 .
\]

Using Bayes' law,
\[
p(d,X) = p(d\mid X)\,p(X) = p(X\mid d)\,p(d),
\]
so that
\[
p(d\mid X) = \frac{p(X\mid d)\,p(d)}{p(X)} .
\]

Given an observation $X$, choose the $d$ that maximizes $p(d\mid X)$. On the
right-hand side:
\begin{itemize}[nosep]
  \item $p(d)=\tfrac12$ is a constant, independent of $d$;
  \item $p(X)$ is independent of $d$.
\end{itemize}
Therefore
\[
p(d\mid X)\ \propto\ p(X\mid d).
\]

Take the optimal $d$ given observed $X$ to be the value $d=d^{*}$ that
maximizes $p(d\mid X)$:
\[
d^{*}
= \argmax_{d}\ p(d\mid X)
\quad\text{(not known directly)}
= \argmax_{d}\ p(X\mid d)
\quad\text{(known directly),}
\]
since $p(X\mid d=P)=p(x)$ and $p(X\mid d=Q)=q(x)$.

%---------------------------------------------------------------
\subsection{The decision rule}
%---------------------------------------------------------------

\[
\text{Given } x,\ \text{choose}\quad
\begin{cases}
P & \text{if } p(x) > q(x),\\[2pt]
Q & \text{if } q(x) > p(x).
\end{cases}
\]

\paragraph{Performance of the decision rule.}
What is $p^{\mathrm{succ}}$ for given $p(x),\,q(x)$? Generally, computing the
performance is more involved than deriving the optimal decision rule. Write
\[
p^{\mathrm{succ}} = P(d^{*}=d_{c}),
\]
where $d_{c}\in\{P,Q\}$ is the chosen source distribution of $X$, and we want
the marginal over $P$ and $Q$.

Write the full joint distribution of all variables, $P(d^{*},x,d_{c})$. Since
$d^{*}$ is a function of $x$ and we are interested in the event $d^{*}=d_{c}$,
\[
p^{\mathrm{succ}} = P(d^{*}=d_{c})
= \sum_{\substack{\forall x\\ d_{c}\in\{P,Q\}}}
P\!\left(d^{*}(x)=d_{c},\,x,\,d_{c}\right).
\]
Using the Markov decomposition $P(A,B,C)=P(A\mid B,C)\,P(B\mid C)\,P(C)$,
\[
p^{\mathrm{succ}}
= \sum_{x,\,d_{c}}
P\!\left(d^{*}(x)=d_{c}\mid x,d_{c}\right)\,
\underbrace{P(x\mid d_{c})}_{\substack{p(x)\ \text{for } d_{c}=P\\ q(x)\ \text{for } d_{c}=Q}}\;
\underbrace{P(d_{c})}_{\tfrac12}.
\]

But we can write
\[
P\!\left(d^{*}(x)=d_{c}\mid x,d_{c}\right)
=
\begin{cases}
\indic\!\left[\,p(x)>q(x)\,\right] & \text{when } d_{c}=P,\\[4pt]
\indic\!\left[\,q(x)>p(x)\,\right] & \text{when } d_{c}=Q,
\end{cases}
\]
where $\indic[\cdot]$ is the indicator function. So
\[
p^{\mathrm{succ}}
= \sum_{x:\,p(x)>q(x)} 1\cdot p(x)\cdot\tfrac12
\;+\;
\sum_{x:\,q(x)>p(x)} 1\cdot q(x)\cdot\tfrac12
= \frac12\!\left(\sum_{x:\,p>q} p(x) + \sum_{x:\,q>p} q(x)\right),
\]
which gives the compact form
\[
\boxed{\; p^{\mathrm{succ}} = \frac12\sum_{x}\max\!\big(p(x),\,q(x)\big). \;}
\]

Using the identity $\max(a,b)=\dfrac{a+b+|a-b|}{2}$, and noting
$\sum_x p(x)=\sum_x q(x)=1$,
\[
p^{\mathrm{succ}}
= \frac14\sum_{x}\Big(p(x)+q(x)+|p(x)-q(x)|\Big)
= \frac12\!\left(1 + \frac12\sum_{x}|p(x)-q(x)|\right).
\]

%===============================================================
\section{Total Variation}
%===============================================================

Define the \emph{total variation} distance as
\[
\boxed{\;
TV(p,q) = \frac12\sum_{x}|p(x)-q(x)|
= \frac12\,\onenorm{p(x)-q(x)},
\;}
\]
where $\onenorm{f} = \sum_{x}|f(x)| = \displaystyle\int|f(x)|\,dx$ is the
one-norm. Then
\[
\boxed{\;
p^{\mathrm{succ}} = \frac12\!\left(1 + \frac12\,TV(p,q)\right).
\;}
\]

%---------------------------------------------------------------
\subsection*{Properties of $TV(p,q)=\tfrac12\onenorm{p-q}$}
%---------------------------------------------------------------

\begin{enumerate}[label=\textbf{\arabic*.}, leftmargin=*]
\item \textbf{Minimum value $0$.}
For each $x$ in the sum, the term is minimized by having $p(x)=q(x)$. Then the
distributions are indistinguishable, and
\[
p^{\mathrm{succ}} = \frac12\!\left(1 + \frac12\cdot 0\right) = \frac12
\quad\rightarrow\ \text{random choice.}
\]

\item \textbf{Maximum value $1$.}
Occurs when the supports of $p(x)$ and $q(x)$ are disjoint, where
$\supp f(x)=\{x\mid f(x)\neq 0\}$. Suppose $p,q$ are disjoint:
\begin{align*}
\frac12\sum_{x}|p(x)-q(x)|
&= \frac12\!\left(
\sum_{\substack{x\\ p\neq 0,\ q=0}}|p(x)-q(x)|
+ \sum_{\substack{x\\ p=0,\ q\neq 0}}|p(x)-q(x)|
\right)\\
&= \frac12\sum_{\substack{x\\ p\neq 0,\ q=0}} p(x)
+ \frac12\sum_{\substack{x\\ p=0,\ q\neq 0}} q(x)
\qquad(\text{disjoint})\\
&= \frac12 + \frac12 = 1.
\end{align*}
(Overlap always decreases the sum.)

\item \textbf{$TV(p,q)=\tfrac12\onenorm{p-q}$ is a true metric distance,} so
the pdf's form a metric space:
\begin{enumerate}[label=(\roman*), nosep]
  \item $TV(p,q)\ge 0$;
  \item $TV(p,q)=TV(q,p)$;
  \item $TV(p,q)=0 \iff p=q$;
  \item triangle inequality.
\end{enumerate}
(One can do limits, completion, etc.\ as usual in abstract metric spaces.)
\end{enumerate}

%===============================================================
\section{Quantum Binary Discrimination}
%===============================================================

Same idea as classical discrimination, but between two quantum states:
density operators / matrices $\rho$ and $\sigma$.

%---------------------------------------------------------------
\subsection{Review}
%---------------------------------------------------------------

An observable is a Hermitian operator / matrix $A$. It admits an
eigendecomposition:
\[
A = \sum_{a} a\,\ket{\psi_a}\!\bra{\psi_a}
  = \sum_{a} a\,P_a,
\]
where the $a$ are real eigenvalues and measured physical outcomes (e.g.\
energy, position), the $\ket{\psi_a}$ are orthonormal eigenstates, and $P_a$
is a projection operator. Sometimes written $\sum_a a\,\ket{a}\!\bra{a}$, with
$a$ now a label (a real number).

An ensemble of quantum states, each with an assigned probability
$\{p_i,\,\ket{\psi_i}\}$, is represented as a density operator / matrix:
\[
\rho = \sum_{i} p_i\,\ket{\psi_i}\!\bra{\psi_i}
\quad\longrightarrow\ \text{a quantum version of a classical pd function.}
\]

Given a quantum state $\rho$ and an observable $A$, the probability of
measuring an outcome $a$ is
\[
P(a\mid\rho) = \Tr P_a\rho.
\]
Why? Because
\begin{align*}
\Tr P_a\rho
&= \Tr\Big(\ket{a}\!\bra{a}\cdot\sum_i p_i\,\ket{\psi_i}\!\bra{\psi_i}\Big)
= \Tr\sum_i p_i\,\braket{a|\psi_i}\braket{\psi_i|a}\\
&= \sum_i p_i\,|\!\braket{a|\psi_i}\!|^{2}
= \sum_i p(\psi_i)\,p(a\mid\psi_i),
\end{align*}
so that
\[
\boxed{\;\Tr P_a\rho = p(a).\;}
\]
So projections like $P_a$ describe these kinds of measurement outcomes.

\paragraph{Properties of the projectors.}
\begin{enumerate}[label=(\arabic*), nosep]
  \item $P_a \ge 0$ \quad (positive semi-definite);
  \item $P_a^{2}=P_a$, \quad $P_a P_b = P_a\,\delta_{ab}$ \quad
        (projections $+$ orthogonality);
  \item $\displaystyle\sum_a P_a = I$ \quad (completeness).
\end{enumerate}

%---------------------------------------------------------------
\subsection{POVMs}
%---------------------------------------------------------------

In more general cases --- coarse graining, noisy measurements, or attaching an
ancilla system, performing a projective measurement, and discarding --- we use
more general operators: POVMs, $E_k$.

\paragraph{Properties.}
\begin{enumerate}[label=(\arabic*), nosep]
  \item Born outcome probability rule: $\ P(k\mid\rho) = \Tr E_k\rho$;
  \item $E_k \ge 0$;
  \item $\displaystyle\sum_k E_k = I$.
\end{enumerate}
(Properties 1--3 are the same as for PVM projective measurements.) The
following differ from PVMs:
\begin{enumerate}[label=(\arabic*), start=4, nosep]
  \item Not orthogonal: $\ E_\ell E_m \neq \delta_{\ell m} E_\ell$ in general;
  \item Not projections: $\ E_k^{2} \neq E_k$ in general.
\end{enumerate}

%===============================================================
\section{Two-Outcome POVM Discrimination}
%===============================================================

We are given two possible states $\rho_0,\rho_1$. We will be handed one to use,
with probability $p_0$ for $\rho_0$ and $p_1$ for $\rho_1$, where $p_0+p_1=1$.
We perform a two-outcome measurement to decide which state, using a POVM
$\{E_0,E_1\}$. Let $p_0=p$ and $p_1=1-p$.

\begin{enumerate}[label=\textbf{\arabic*.}, leftmargin=*]
\item We have $E_0+E_1=I$ and $E_0,E_1\ge 0$. So the POVM is $\{E_0,\,I-E_0\}$;
call $E_0=E$. Then $I-E\ge 0 \Rightarrow E\le I$, so
\[
0 \le E \le I \quad\longrightarrow\quad
\text{all eigenvalues of } E \text{ satisfy } 0\le\lambda\le 1.
\]

\item We have two quantum states prepared: $\rho_0$ with $P(\rho_0)=p$, and
$\rho_1$ with $P(\rho_1)=1-p$. (Compare the classical case of pdf's $\{p,q\}$
with probability $\tfrac12$ each.)

\item Perform a measurement:
\[
\begin{aligned}
\text{Outcome } E_0=E \ &\Rightarrow\ \text{decide } \rho_0 \text{ is the quantum state},\\
\text{Outcome } E_1=I-E \ &\Rightarrow\ \text{decide } \rho_1 \text{ is the quantum state}.
\end{aligned}
\]

\item We have two ways of deciding correctly:
\[
\begin{aligned}
&\text{True state } \rho_0 \ (\text{prob. } p):
&&\text{Prob.\ of outcome } E_0=E \text{ is } \Tr E\rho_0;\\
&\text{True state } \rho_1 \ (\text{prob. } 1-p):
&&\text{Prob.\ of outcome } E_1=I-E \text{ is } \Tr(I-E)\rho_1.
\end{aligned}
\]
Total:
\[
\boxed{\; P^{\mathrm{corr}} = p\,\Tr E\rho_0 + (1-p)\,\Tr(I-E)\rho_1. \;}
\]
We want to choose an $E$, given $\{p,\rho_0,\rho_1\}$, that maximizes
$P^{\mathrm{corr}}$ --- an optimization problem.

\item Now expand and rearrange; we get something similar to the classical case:
\begin{align*}
P^{\mathrm{corr}}
&= p\,\Tr E\rho_0 + (1-p)\,\Tr(I-E)\rho_1\\
&= p\,\Tr E\rho_0 + (1-p)\big[\Tr\rho_1 - \Tr E\rho_1\big]\\
&= p\,\Tr E\rho_0 + \Tr\rho_1 - \Tr E\rho_1 - p\,\Tr\rho_1 + p\,\Tr E\rho_1\\
&= p\,\Tr E\rho_0 + (1-p) - (1-p)\,\Tr E\rho_1\\
&= (1-p) + \Tr E\big(p\rho_0 - (1-p)\rho_1\big).
\end{align*}
So
\[
\boxed{\; P^{\mathrm{corr}} = (1-p) + \Tr E\big(p_0\rho_0 - p_1\rho_1\big) \;}
\xrightarrow{\ p=1/2\ }
\frac12 + \frac12\,\Tr E(\rho_0-\rho_1),
\]
compare with the classical
$\ \tfrac12 + \tfrac12\sum_x|p(x)-q(x)| = \tfrac12 + \tfrac12\onenorm{p-q}$.
\end{enumerate}

Define the \emph{Helstrom operator}
\[
\boxed{\; \Delta = p_0\rho_0 - (1-p)\rho_1 = p_0\rho_0 - p_1\rho_1, \;}
\]
which encodes the \emph{known} problem information. So
\[
\boxed{\; P^{\mathrm{corr}} = 1-p + \Tr E\Delta \;}
\quad\longrightarrow\quad
\text{for the given }\Delta,\ \text{adjust } E \text{ with } 0\le E\le I
\text{ to maximize } P^{\mathrm{corr}}.
\]
Note we are optimizing over an operator (matrix) in an expression involving a
trace, under the psd constraints $0\le E\le I$.

%===============================================================
\section{The Optimization Problem}
%===============================================================

For given $\Delta = p_0\rho_0 - p_1\rho_1$, find
\[
\boxed{\; E^{*} = \argmax_{0\le E\le I}\ \Tr E\Delta, \;}
\]
where $0\le E\le I$ means all eigenvalues of $E$ lie in $[0,1]$.

We know $\rho_0,\rho_1$ are psd (all eigenvalues $\lambda\ge 0$) and
$\Tr\rho_0=\Tr\rho_1=1$. But the weighted difference $p_0\rho_0-p_1\rho_1$ is
still Hermitian, yet not necessarily psd (some eigenvalues $<0$ are possible).

Suppose we set $E=I$. Then
\[
\Tr I\Delta = \Tr\Delta = \sum_i \delta_i
\qquad(\delta_i = \text{eigenvalues of }\Delta).
\]
Generally some $\delta_i>0$ and some $\delta_i<0$, so the sum will be at least
partially decreased by the negative $\delta_i$.

\begin{quote}
\textbf{Main idea:} choose $E$ as a \emph{projector onto only the positive
eigenvalue subspace} of $\Delta$. Then $\Tr E\Delta$ will sum only positive
values. Design $E$ to maximize this.
\end{quote}

%---------------------------------------------------------------
\subsection{The optimal $E$}
%---------------------------------------------------------------

\paragraph{Claim.}
Eigendecompose $\Delta$:
\[
\Delta = \sum_i \delta_i\,\ket{i}\!\bra{i}
= \sum_j \delta_j^{+}\,\ket{j^{+}}\!\bra{j^{+}}
+ \sum_k \delta_k^{-}\,\ket{k^{-}}\!\bra{k^{-}},
\]
where $\{\delta_j^{+}\}$ are the positive eigenvalues and $\{\delta_k^{-}\}$
are the zero or negative eigenvalues. So $\Delta=\Delta_{+}+\Delta_{-}$ with
\[
\Delta_{+} = \sum_j \delta_j^{+}\,\ket{j^{+}}\!\bra{j^{+}},
\qquad
\Delta_{-} = \sum_k \delta_k^{-}\,\ket{k^{-}}\!\bra{k^{-}}.
\]
Then
\[
\boxed{\; E^{*} = \sum_j \ket{j^{+}}\!\bra{j^{+}} \equiv \Pi_{+}, \;}
\]
a projection onto the positive eigenspace of $\Delta$.

\paragraph{Why.}
\begin{enumerate}[label=\textbf{\arabic*.}, leftmargin=*]
\item \emph{Notation.}
\begin{enumerate}[label=(\roman*), nosep]
  \item $A>0$ means, for all unit $v$, $\ \braket{v|A|v}>0 \quad \forall v$.
  \item $A<I$ means, for all $v$ with $|v|=1$,
  \[
  A-I<0 \ \Longrightarrow\ \braket{v|A-I|v}<0
  \ \Longrightarrow\ \braket{v|A|v} - \underbrace{\braket{v|v}}_{1} < 0
  \ \Longrightarrow\ \braket{v|A|v}<1 \quad\forall v.
  \]
\end{enumerate}

\item \emph{Upper bound.}
We have $\Delta=\sum_i\delta_i\ket{i}\!\bra{i}$, so
\[
\Tr E\Delta
= \Tr E\Big(\sum_i\delta_i\ket{i}\!\bra{i}\Big)
= \sum_i \delta_i\,\Tr\big(E\ket{i}\!\bra{i}\big)
= \sum_i \delta_i\,\braket{i|E|i}.
\]
We know (i) $0\le\braket{i|E|i}\le 1$ and (ii) some $\delta_i>0$, some
$\delta_i\le 0$. We can maximize this sum term by term by setting
\[
\braket{i|E|i}=1 \ \text{for } \delta_i>0,
\qquad
\braket{i|E|i}=0 \ \text{for } \delta_i\le 0.
\]
(We don't know yet that this is achievable, but it does set an upper bound.)
Therefore
\[
\boxed{\; \Tr E\Delta \le \sum_i \delta_i^{+} = \Tr\Delta_{+}. \;}
\]

\item \emph{Achievability.}
This bound is achieved by setting
\[
\boxed{\; E^{*} = \sum_j \ket{j^{+}}\!\bra{j^{+}} \equiv \Pi_{+}. \;}
\]
Why? Because
\begin{align*}
\Tr E^{*}\Delta
&= \Tr\Big(\sum_j \ket{j^{+}}\!\bra{j^{+}}\cdot\sum_i\delta_i\ket{i}\!\bra{i}\Big)
= \Tr\sum_{j,i}\delta_i\,\ket{j^{+}}\!\underbrace{\braket{j^{+}|i}}_{\delta_{j^{+}i}}\!\bra{i}\\
&= \sum_i \delta_i^{+} = \Tr\Delta_{+}.
\end{align*}
Hence the optimal POVM element is
\[
\boxed{\; E^{*} = \sum_i \ket{i^{+}}\!\bra{i^{+}}. \;}
\]
\end{enumerate}

%===============================================================
\section{The Probability of Success}
%===============================================================

Now look at $p^{\mathrm{succ}}$. For $\Delta = p_0\rho_0 - p_1\rho_1$,
\[
p^{\mathrm{succ}} = 1-p + \Tr E^{*}\Delta = 1-p + \Tr\Delta_{+}.
\]
Now
\[
\Tr\Delta_{+} = \sum_{i:\,\delta_i>0}\delta_i
= \frac12\!\left(\sum_i|\delta_i| + \sum_i\delta_i\right)
= \frac12\big(\onenorm{\Delta} + \Tr\Delta\big),
\]
and
\[
\Tr\Delta = \Tr\big(p\rho_0 - (1-p)\rho_1\big)
= p\,\Tr\rho_0 - (1-p)\,\Tr\rho_1
= p - (1-p) = 2p-1,
\]
so that
\[
\Tr\Delta_{+} = \frac12\big(\onenorm{\Delta} + 2p-1\big).
\]
Therefore
\[
p^{\mathrm{succ}}
= 1-p + \frac12\onenorm{\Delta} + p - \frac12
= \frac12\big(1 + \onenorm{\Delta}\big).
\]
So the optimal probability of success in binary discrimination is
\[
\boxed{\;
p^{\mathrm{succ}} = \frac12\Big(1 + \onenorm{\,p\rho_0 + (1-p)\rho_1\,}\Big).
\;}
\]

\paragraph{Comparison to the classical $p=\tfrac12$ case.}
\[
\boxed{\;
p^{\mathrm{succ}} = \frac12\!\left(1 + \frac12\onenorm{p-q}\right)
\;}
\qquad\text{(classical),}
\]
\[
\boxed{\;
p^{\mathrm{succ}} = \frac12\!\left(1 + \frac12\onenorm{\rho_0-\rho_1}\right)
\;}
\qquad\text{(quantum).}
\]

%===============================================================
\section{Quantum Trace Distance of States}
%===============================================================

(For $p=\tfrac12$.) Define
\[
T(\rho,\sigma) = \frac12\onenorm{\rho-\sigma},
\qquad \onenorm{A}\equiv\sum_i|\lambda_i|,
\]
alongside the classical
\[
TV(p,q) = \frac12\onenorm{p-q},
\qquad \onenorm{f}\equiv\sum_x|f(x)|.
\]

\paragraph{Same general properties.}
\begin{enumerate}[label=\textbf{\arabic*.}, leftmargin=*]
  \item Both are proper metric distances.
  \item Minimum $0$: $\ \rho=\sigma \iff T=0 \ \rightarrow\ p^{\mathrm{succ}}=\tfrac12$.
  \item Maximum $1$: $\ \rho,\sigma$ perfectly distinguishable with
        $p^{\mathrm{succ}}=1$.
  \item Perfect distinguishability also has a disjointness requirement.
\end{enumerate}

\paragraph{Perfect distinguishability.}
One can show that for $T(\rho,\sigma)=1$:
\[
\begin{aligned}
&\text{(1) for } \delta_i>0:\ \braket{i|\sigma|i}=0,\\
&\text{(2) for } \delta_i<0:\ \braket{i|\rho|i}=0.
\end{aligned}
\]
Or, equivalently:
\begin{enumerate}[label=(\arabic*), nosep]
  \item all positive eigenvectors of $\rho-\sigma$ belong to the image of $\rho$;
  \item all negative eigenvectors of $\rho-\sigma$ belong to the image of $\sigma$.
\end{enumerate}

\end{document}
